% Part: applied-modal-logics
% Chapter: temporal-logic
% Section: language-epistemic-logic

\documentclass[../../../include/open-logic-section]{subfiles}

\begin{document}

\olfileid{aml}{tl}{sem}

\olsection{时态逻辑的语义}

\begin{defn}
时态逻辑的基本语言包含以下要素：
\begin{enumerate}
  \tagitem{prvFalse}{表示假的命题常项~$\lfalse$。}{}
  \tagitem{prvTrue}{表示真的命题常项~$\ltrue$。}{}
  \item 一个可数无穷的命题变元集合：$\Obj p_0$、$\Obj p_1$、$\Obj p_2$……
  \item 命题联结词：\startycommalist
  \iftag{prvNot}{\ycomma $\lnot$（否定）}{}%
  \iftag{prvAnd}{\ycomma $\land$（合取）}{}%
  \iftag{prvOr}{\ycomma $\lor$（析取）}{}%
  \iftag{prvIf}{\ycomma $\lif$（条件句）}{}%
  \iftag{prvIff}{\ycomma $\liff$（双条件句）}{}。
  \item 过去时态算子 $\Ptemp$（曾有）和 $\Htemp$（过去一直）。
  \item 未来时态算子 $\Ftemp$（将有）和 $\Gtemp$（将来一直）。
\end{enumerate}
\end{defn}

稍后，我们将讨论增加其他种类模态算子的可能性。

\begin{defn}
时态语言的\emph{公式}归纳定义如下：
\begin{enumerate}
\tagitem{prvFalse}{$\lfalse$ 是一个原子公式。}{}

\tagitem{prvTrue}{$\ltrue$ 是一个原子公式。}{}

\item 每个命题变元 $\Obj p_i$ 都是一个（原子）公式。

\tagitem{prvNot}{若 $!A$ 是公式，则 $\lnot !A$ 是公式。}{}

\tagitem{prvAnd}{若 $!A$ 和 $!B$ 是公式，则 $(!A \land !B)$ 是公式。}{}

\tagitem{prvOr}{若 $!A$ 和 $!B$ 是公式，则 $(!A \lor !B)$ 是公式。}{}

\tagitem{prvIf}{若 $!A$ 和 $!B$ 是公式，则 $(!A \lif !B)$ 是公式。}{}

\tagitem{prvIff}{若 $!A$ 和 $!B$ 是公式，则 $(!A \liff !B)$ 是公式。}{}

\item 若 $!A$ 是公式，则 $\Ptemp !A$、$\Htemp !A$、$\Ftemp !A$、$\Gtemp !A$ 均为公式。

\tagitem{limitClause}{其他均非公式。}{}
\end{enumerate}
\end{defn}

与其他内涵逻辑一样，时态逻辑的语义也是借助关系模型给出的。

\begin{defn}
  时态语言的一个\emph{模型}是一个三元组 $\mModel{M} = \tuple{T, \prec, V}$，其中
  \begin{enumerate}
  \item $T$ 是一个非空集合，解释为时间点的集合。
  \item $\prec$ 是 $T$ 上的一个二元关系。
  \item $V$ 是一个函数，对每个命题变元~$p$ 指派一个时间点子集 $V(p)$。
  \end{enumerate}
  当 $t \prec t'$ 成立时，我们称 $t$ \emph{前驱于}~$t'$。
  当 $t \in V(p)$ 时，我们称 $p$ \emph{在}~$t$ \emph{处为真}。
\end{defn}

目前，你会注意到我们没有对前驱关系 $\prec$ 施加任何条件。这意味着现阶段我们的时态模型结构没有任何限制，因此我们可以有时间为线性、分支、循环或具有任何结构的模型。

与正规模态逻辑一样，每个时态模型决定了其中哪些公式在哪些点处为真。对于关系语义的基本概念，我们沿用相同的记法“模型 $\mModel{M}$ 使公式~$!A$ 在点~$t$ 处为真”。该关系是归纳定义的，并且对所有非模态算子而言，与正规模态逻辑的情形完全相同。

\begin{defn}\ollabel{defn:tmodels}
  公式~$!A$ 在模型~$\mModel M$ 中的点~$t$ 处的\emph{真}，记作 $\mSat{M}{!A}[t]$，归纳定义如下：
  \begin{enumerate}
  \tagitem{prvFalse}{\indcase{!A}{\lfalse}{永不满足 $\mSat{M}{\lfalse}[t]$}。}{}
  \tagitem{prvTrue}{\indcase{!A}{\ltrue}{总是满足 $\mSat{M}{\ltrue}[t]$}。}{}
  \item $\mSat{M}{p}[t]$ 当且仅当 $t \in V(p)$。
  \tagitem{prvNot}{\indcase{!A}{\lnot !B}{$\mSat{M}{\indfrm}[t]$ 当且仅当 $\mSat/{M}{!B}[t]$}。}{}
  \tagitem{prvAnd}{\indcase{!A}{(!B \land !C)}{$\mSat{M}{\indfrm}[t]$ 当且仅当 $\mSat{M}{!B}[t]$ 且 $\mSat{M}{!C}[t]$}。}{}
  \tagitem{prvOr}{\indcase{!A}{(!B \lor !C)}{$\mSat{M}{\indfrm}[t]$ 当且仅当 $\mSat{M}{!B}[t]$ 或 $\mSat{M}{!C}[t]$（或两者）}。}{}
  \tagitem{prvIf}{\indcase{!A}{(!B \lif !C)}{$\mSat{M}{\indfrm}[t]$ 当且仅当 $\mSat/{M}{!B}[t]$ 或 $\mSat{M}{!C}[t]$}。}{}
  \tagitem{prvIff}{\indcase{!A}{(!B \liff !C)}{$\mSat{M}{\indfrm}[t]$ 当且仅当要么 $\mSat{M}{!B}[t]$ 与 $\mSat{M}{!C}[t]$ 同时成立，要么 $\mSat{M}{!B}[t]$ 与 $\mSat{M}{!C}[t]$ 均不成立}。}{}
  \item\ollabel{defn:sub:mmodels-p}
    \indcase{!A}{\Ptemp !B}{$\mSat{M}{\indfrm}[t]$ 当且仅当对某个满足 $t' \prec t$ 的 $t' \in T$，有 $\mSat{M}{!B}[t']$}。
\item\ollabel{defn:sub:mmodels-h}
    \indcase{!A}{\Htemp !B}{$\mSat{M}{\indfrm}[t]$ 当且仅当对所有满足 $t' \prec t$ 的 $t' \in T$，都有 $\mSat{M}{!B}[t']$}。
   \item\ollabel{defn:sub:mmodels-f}
    \indcase{!A}{\Ftemp !B}{$\mSat{M}{\indfrm}[t]$ 当且仅当对某个满足 $t \prec t'$ 的 $t' \in T$，有 $\mSat{M}{!B}[t']$}。
\item\ollabel{defn:sub:mmodels-g}
    \indcase{!A}{\Gtemp !B}{$\mSat{M}{\indfrm}[t]$ 当且仅当对所有满足 $t \prec t'$ 的 $t' \in T$，都有 $\mSat{M}{!B}[t']$}。
  \end{enumerate}
\end{defn}

根据语义，或许能看出算子 $\Ptemp$ 与~$\Htemp$ 互为对偶，算子 $\Ftemp$ 与 $\Gtemp$ 也互为对偶，从而我们可以将 $\Htemp  !A$ 定义为 $\lnot \Ptemp \lnot !A$，$\Gtemp$ 与~$\Ftemp$ 之间的关系也是如此。

\end{document}